\documentclass[12pt]{article}
\usepackage{pst-all}
\usepackage[french]{babel}
\usepackage{tikz}
\usepackage{tkz-tab}
\usepackage[T1]{fontenc}
\usepackage{eurosym}
\usepackage{amssymb}
\usepackage{float}
\usepackage{blkarray}
\usepackage{dsfont}
\usepackage{amsfonts}
\usepackage{amsmath}
\usepackage{eurosym}
\usepackage{array}
\usepackage{graphicx}
\parindent=0cm
\usepackage{geometry}
\usepackage{color}
\usepackage{graphicx}
\geometry{hmargin=1.5cm,vmargin=1.5cm}
\usepackage{amsmath, amssymb, amsfonts, amscd, amsthm}
\usepackage{pifont}
\usepackage{fancyhdr} 
\pagestyle{fancy}
\usepackage{multirow}
\renewcommand{\headrulewidth}{1pt}
\renewcommand{\footrulewidth}{1pt}
\fancyhead[L]{~}
\fancyhead[R]{~}
\fancyhead[C]{$5^e$}





\usepackage[utf8]{inputenc}
\usepackage{tcolorbox}% charge aussi pgf, verbatim, etoolbox et environ
\tcbuselibrary{theorems}% charge aussi amsmath
\usepackage{xcolor}

\definecolor{propcolor}{rgb}{.804,.064,.464}
\definecolor{defcolor}{rgb}{0,.545,0}
\definecolor{thcolor}{rgb}{0.932,0,0}
\definecolor{remcolor}{rgb}{1,0.5,0}
\definecolor{excolor}{rgb}{.094,.455,.804}

\usepackage{cleveref}

\newtcbtheorem[crefname={definition}{definitions}]{Definition}{Definition}{defstyle}{def}


\usepackage{varwidth} 
\tcbuselibrary{skins}

\newtcbtheorem[use counter from=Definition]{defin}{Définition}
{enhanced,frame empty, theorem name,interior empty,colframe=defcolor,coltitle=defcolor,fonttitle=\bfseries,
	colbacktitle=white,
	borderline={0.5mm}{0mm}{defcolor},
	%borderline={0.5mm}{0mm}{ForestGreen!50!white,dashed},
	attach boxed title to top left={yshift=-2mm, xshift = 2mm},
	boxed title style={boxrule=0.4pt}, varwidth boxed title}{theo}

\newtcbtheorem[use counter from=Definition]{theo}{Théorème}
{enhanced,frame empty,theorem name,interior empty,colframe=thcolor,coltitle=thcolor,fonttitle=\bfseries,
	colbacktitle=white,
	borderline={0.5mm}{0mm}{thcolor},
	%borderline={0.5mm}{0mm}{ForestGreen!50!white,dashed},
	attach boxed title to top left={yshift=-2mm, xshift = 2mm},
	boxed title style={boxrule=0.4pt}, varwidth boxed title}{theo}

\newtcbtheorem[use counter from=Definition]{prop}{Propriété}
{enhanced,frame empty,theorem name,interior empty,colframe=propcolor,coltitle=propcolor,fonttitle=\bfseries,
	colbacktitle=white,
	borderline={0.5mm}{0mm}{propcolor},
	%borderline={0.5mm}{0mm}{ForestGreen!50!white,dashed},
	attach boxed title to top left={yshift=-2mm, xshift = 2mm},
	boxed title style={boxrule=0.4pt}, varwidth boxed title}{theo}


\newtcbtheorem[use counter from=Definition]{ex}{Exemple}
{enhanced,frame empty,theorem name,interior empty,colframe=excolor,coltitle=excolor,fonttitle=\bfseries,
	colbacktitle=white,
	borderline={0.5mm}{0mm}{excolor},
	%borderline={0.5mm}{0mm}{ForestGreen!50!white,dashed},
	attach boxed title to top left={yshift=-2mm, xshift = 2mm},
	boxed title style={boxrule=0.4pt}, varwidth boxed title}{theo}

\newtcbtheorem[use counter from=Definition]{rem}{Remarque}
{enhanced,frame empty,theorem name,interior empty,colframe=remcolor,coltitle=remcolor,fonttitle=\bfseries,
	colbacktitle=white,
	borderline={0.5mm}{0mm}{remcolor},
	%borderline={0.5mm}{0mm}{ForestGreen!50!white,dashed},
	attach boxed title to top left={yshift=-2mm, xshift = 2mm},
	boxed title style={boxrule=0.4pt}, varwidth boxed title}{theo}


\newif\iftrou\troutrue
%mettre true à la place de false pour mettre des pointillés

\newlength{\trouwidth}

\newcommand{\trou}[1]{%
	\iftrou%
	\settowidth{\trouwidth}{#1}%
	\hbox to \trouwidth{\dotfill}%
	\else%
	#1%
	\fi}


\newcommand\pointille[1][1cm]{\makebox[#1]{\dotfill}}


\begin{document}


%%%%%%%%%%%%%%%%%%%%%%%%%%%%%%%%%%%%%%%%%%%%%%%%%%%%%%%%%%%%


\begin{center}
		\rule{\linewidth}{0.5mm} \\[0.4cm]
		{ \huge \bfseries Probabilités\\[0.4cm] }
		\rule{\linewidth}{0.5mm} \\
\end{center}

%%%%%%%%%%%%%%%%%%%%%%%%%%%%%%%%%%%%%%%%%%%%%%%%%%%%%%%%%%%%%%%%%%%%
%%%%%%%%%%%%%%%%%%%%%%%%%%%%%%%%%%%%%%%%%%%%%%%%%%%%%%%%%%%%%%%%%%%%


\section{Expérience aléatoire}

\begin{defin}{}{}
	Une \textbf{expérience aléatoire} est une expérience dans laquelle intervient le hasard: on ne peut pas prévoir le résultat à l'avance.\\
	Les différents résultats possibles d'une expérience aléatoire sont appelés des \textbf{issues}.
\end{defin}

\begin{ex}{}{}
	\begin{enumerate}
		\item On lance une pièce de monnaie et on observe la face du dessus. Il y a deux issues possibles : \trou{\textcolor{brown}{Pile}} et \trou{\textcolor{purple}{Face}}. On ne peut pas savoir à l'avance laquelle des deux on va obtenir. 
		\item On lance un dé de $6$ faces et on observe le numéro inscrit sur la face du dessus. Il y a \trou{$6$} issues possibles: \trou{$1$, $2$, $3$, $4$, $5$ et $6$}. On ne peut pas savoir à l'avance laquelle de ces $6$ issues on va obtenir.
	\end{enumerate}
\end{ex}

\begin{defin}{}{}
	Un \textbf{événement} est une condition qui, selon l'issue de l'expérience aléatoire, est \textbf{réalisé} ou n'est \textbf{pas réalisé}.\\
		Une issue peut également être considérée comme un \textbf{événement}.
\end{defin}


\begin{ex}{}{}
		Lors d'un lancer de dé, on s'intéresse à la question suivante : "le résultat obtenu est-il pair ?"\\
		On parle alors de l'événement "Obtenir un résultat pair", c'est-à-dire "Obtenir \trou{$2$, $4$ ou $6$}".\\
		Si on obtient l'une des issues $2$, $4$ ou $6$, on dit que l'\textbf{événement} "Obtenir un résultat pair" est \textbf{réalisé}.\\
		Si on obtient l'une des issues $1$, $3$ ou $5$, on dit que l'\textbf{événement} "Obtenir un nombre pair" n'est \textbf{pas réalisé}.
\end{ex}


\begin{ex}{}{}
	On tire une fléchette les yeux bandés vers la cible ci-dessous et on observe le nombre de points obtenus.
	Si la fléchette ne se plante pas sur la cible, on n'obtient aucun point. Les issues de cette expérience aléatoire sont : \trou{$0$, $20$, $40$, $50$, $80$ et $100$}.
	\begin{center}
		\includegraphics[width=0.2\linewidth]{flechette}
		\end{center}
	On s'intéresse à l'événement "Obtenir $40$ points" : si on obtient l'issue $40$ lors de cette expérience, cet événement sera réalisé. Si on obtient une autre \trou{issue}, cet événement ne sera pas réalisé. 
\end{ex}




%%%%%%%%%%%%%%%%%%%%%%%%%%%%%%%%%%%%%%%%%%%


\section{Probabilité d'un événement}

\begin{defin}{}{}
	La \textbf{probabilité} d'un événement peut s'interpréter comme la "proportion de chances" que cet événement se réalise. C'est un nombre \textbf{compris entre $0$ et $1$}.\\
	Plus un événement a de chances de se réaliser, plus sa probabilité est proche de $1$.\\
	Moins il a de chances de se réaliser, plus sa probabilité est proche de $0$.
		\begin{center}
		\includegraphics[width=0.6\linewidth]{proba1}
	\end{center}
\end{defin}


\begin{rem}{}{}
	On peut exprimer une probabilité sous plusieurs formes: un nombre décimal, une fraction, un pourcentage, etc.
\end{rem}


\begin{defin}{}{}
	Pour une expérience aléatoire, lorsqu'on a autant de chances d'obtenir une issue ou une autre, on dit qu'il s'agit d'une \textbf{situation d'équiprobabilité}.
\end{defin}

\begin{defin}{}{}
	Si $A$ désigne un événement, on note $\mathbb{P}(A)$ la probabilité de cet événement.
\end{defin}


\begin{ex}{Lancer de dé}{}
	Lorsqu'on lance un dé à 6 faces qui n'est pas truqué, il a autant de chances de tomber sur une face que sur une autre. On a donc $1$ chance sur $6$ d'obtenir le $1$, $1$ chance sur $6$ d'obtenir le $2$, $1$ chance sur $6$ d'obtenir le $3$, etc.\\
	Toutes les issues ont les mêmes chances d'être obtenues. On dit aussi que les issues ont la même probabilité. Cette probabilité est égale à $\displaystyle\frac{1}{6}$.
	\begin{center}
		\includegraphics[width=0.1\linewidth]{proba2}
	\end{center}
\end{ex}

\begin{ex}{Lancer de pièces}{}
	Lorsqu'on lance une pièce de monnaie qui EST truquée, elle n'a pas autant de chances de tomber sur une face que sur l'autre. On a peut-être $40\%$ de chances d'obtenir "Face" et $60\%$ de chances d'obtenir "Pile". Ce n'est pas une situation d'équiprobabilité.
	\begin{center}
		\includegraphics[width=0.2\linewidth]{proba3}
	\end{center}
\end{ex}

\begin{prop}{}{}
	On considère une situation d'\textbf{équiprobabilité} admettant $N$ issues. La probabilité d'un événement est égale au quotient du nombre d'issues qui réalisent cet événement par le nombre $N$.
\end{prop}

\begin{ex}{}{}
	Une urne contient $4$ boules indiscernables au toucher.
		\begin{center}
		\includegraphics[width=0.15\linewidth]{proba4}
	\end{center}
	On tire au hasard une boule de l'urne. \\
	Les boules étant indiscernables au toucher, chaque boule a la même chance d'être tirée, il s'agit donc d'une situation d'équiprobabilité.\\
	Cette expérience admet \textcolor{red}{$4$} issues, on a donc:\\ $\mathbb{P}$(obtenir $1$) =  $\mathbb{P}$(obtenir $2$) =  $\mathbb{P}$(obtenir $3$) =  $\mathbb{P}$(obtenir $4$) = $\displaystyle\frac{1}{\textcolor{red}{4}}$.\\
	Si on s'intéresse maintenant à l'événement "\textcolor{blue}{Obtenir une boule bleue}", on a :\\ 
	$\mathbb{P}$(Obtenir une boule bleue) = $\displaystyle\frac{\textcolor{blue}{2}}{\textcolor{red}{4}}$
\end{ex}




%%%%%%%%%%%%%%%%%%%%%%%%%%%%%%%%%%%%%%%%%%%%%%%%%

\section{Carte mentale}

	\begin{center}
	\includegraphics[width=0.99\linewidth]{cmproba}
\end{center}







%%%%%%%%%%%%%%%%%%%%%%%%%%%%%%%%%%%%%%%%%%%%%%%%%%%%%%%%%%%%%%%%%%%%%%%%%%%%
\end{document}
